% p.67 — الباب التاسع والعشرون في الاصطلاحات والتركيبات (تتمة: التركيبات)
\mspage{67}
% Layout: four columns = two pairs (label | text), right pair first; continues the
% تركيب examples (34)–(48).  Circled example numbers written as (n); the circle of
% (41) is smudged.  Printed page number ۶۷ at top left.  The lower part of the left
% pair is blank.  Layout: the two pairs are set side by side (right pair first) in one
% `tables` row, as in the lithograph.
\begin{tables}
\begin{sarf}{|r|r|}
 & (34) رَأَيْتُ زَيْدًا نَفْسَهُ\\
تأكيد & نَفْسَهُ تأكيد\\
مُؤَكَّد & زَيْدًا مُؤَكَّد\\
 & (35) مَرَرْتُ بِزَيْدٍ وَ عَمْرٍو\\
حرف عطف & واو حرف عطف\\
معطوف به & عَمْرٍو معطوف به\\
معطوف عليه & زيد معطوف عليه\\
 & (36) هٰذَا أَبُوْ حَفْصٍ عُمَرُ\\
مُبَيَّن & ابوحفص مبين\\
عطف بيان & عُمَرُ عطف بيان\\
 & (37) أَكْرَمْتُ زَيْدًا صَدِيْقَكَ\\
مُبَدَّل منه & زيدا مبدل منه\\
بدل & صديق بدل\\
 & (38) أَكَلْتُ الرَّغِيْفَ نِصْفَهُ\\
مبدل منه & الرَّغِيْفَ مبدل مِنْهُ\\
بدل & نِصْفَهُ بَدَلٌ\\
 & (39) زَيْدٌ هُوَ قَائِمٌ\\
ضمير & هُوَ ضمير\\
مرجع ضمير & زيد مرجع ضمير\\
 & (40) أَنَا عَالِمٌ هُوَ جَاهِلٌ\\
ضمير ظاهر & أَنَا . هُوَ . أَنْتَ ضمير ظاهر\\
\end{sarf}
\begin{sarf}{|r|r|}
 & (41) زَيْدٌ ضَرَبَ عَمْرًا\\
ضمير مُسْتَتِر & هُوَ فى ضرب ضمير مُسْتَتِر\\
 & (42) ضَرَبْتَ عَمْرًا\\
ضمير بَارِز & أَنْتَ فى ضَرَبْتَ ضمير بَارِز\\
 & (43) اِنَّهُ زَيْدٌ قَائِمٌ\\
ضمير شأن & ه ضمير شأن\\
 & (44) اِنَّهَا زَيْنَبُ قَائِمَةٌ\\
ضمير قِصَّة & ها ضمير قِصَّة\\
 & (45) رُبَّهُ رَجُلًا\\
ضمير مبهم & ه ضمير مبهم\\
 & (46) نِعْمَ رَجُلًا زَيْدٌ\\
ضمير مبهم & هُوَ فى نِعْمَ ضمير مبهم\\
 & (47) كَانَ زَيْدٌ هُوَ قَائِمًا\\
ضمير فصل & هو ضمير فصل\\
 & (48) كَانَ الْمَطَرُ\\
فعل تام & كان فعل تام\\
فاعل تام & المطر فاعل تام\\
\end{sarf}
\end{tables}
